\subsection{Connections to observed particle and field structure}

The finite constructions reproduce several characteristic mathematical
structures used in descriptions of observed particles and fields. Each
comparison below states the generated record, its calibration and its
physical domain. These results use declared source members; combining
their conclusions into one interacting realization requires the
corresponding source-transfer statements. A structural agreement, an
empirically anchored energy and a held-out numerical prediction are
different tests.

\paragraph{Quantized electric charge.}
The native Fock charge operator has an exact discrete spectrum, including
the charged-lepton unit and quark charges $2/3$ and $-1/3$. Actual charged
arrivals carry the corresponding central current and regional Gauss
flux, with complete failure and reference outcomes. Charge cannot be
continuously divided within a fixed native representation. This is a
finite compact-current realization of charge quantization. It does not
equate texture winding with electric current or by itself identify the
observed electromagnetic field. The complete source calculations are in
\path{model/openwave_charge_protection} and
\path{model/openwave_native_charge_flux}.

\paragraph{Charged rest energy.}
The exact native CAR and local Gauss isometry attach a positive energy
to a dressed charged state, including the root-reference field.
Equal-depth tree dressing adds its explicit electric energy without
changing the matter current on its finite-orbital member. Separately,
the full spatial/photon charged block, without tree dressing, has a
trap-release threshold $16$, conditionally
$0.000509529040252\,\mathrm{GeV}$ under its retained electron conversion.
The empirical electron anchor and its source conventions are explicit;
this energy is not an independent precision electron-mass prediction.
The threshold has no normalizable minimizing state. The distinction
between a free threshold and a localized trapped state is proved in
\path{model/openwave_spatial_completion}. Its enlarged spatial domain
does not inherit the finite-orbital tree energy or neutral-sector bounds.

\paragraph{Particle persistence.}
The complete finite charged packet has a proved persistence domain.
For the spatial/photon source
$H_\Omega=\sqrt{16^2+B^2}+\Omega N_{\rm m}+N_\gamma$,
positivity and oscillator coercivity isolate a spin-half band when
$\min(\Omega,1)>3/64+(3/8192)(5/7)^{5/2}$.
Shell convergence controls the full source in this trap domain.
This is a quantitative stable-band construction. The result supplies
neither a self-bound localized electron after trap removal nor a
nonlinear classical soliton.

\paragraph{Magnetic response near two.}
The actual differentiated current, including the photon field and its
complement, gives a static coefficient in
$[1.99827530667,1.99827540121]$. Global divided-difference estimates and
finite moment witnesses control the infinite-photon source; finite-time
precession uses the same current. The proximity to two is a quantitative
structural comparison with a spin-half magnetic response. Its deviation
does not reproduce the observed electron anomaly. Field normalization,
trap and regulator are declared inputs, as specified in
\path{model/openwave_infinite_current}.

\paragraph{Spin one half.}
The finite-orbital dressed band carries $J^2=3/4$ and the two spin
projections $\pm1/2$ in units of $\hbar$. The native CAR charge map and
longitudinal Gauss isometry preserve this representation on that member.
Independently, the spatial/photon charged block includes complete matter
shells and photon angular momentum, with a spin-half band and full-space
convergence in its stated trap domain. The tree dressing and magnetic
response bounds are not transferred to this enlarged spatial source.
These exact representations agree with the spin structure of a fermionic
particle; physical species identification is an additional comparison.

\paragraph{The double-cover rotation law.}
Literal proper cube rotations act on the kinetic/pairing source and its
filled reference. An odd state changes sign under $2\pi$ relative to an
idle reference and returns under $4\pi$. Complete oriented spin-frame
records return under $2\pi$, while a declared two-copy apolar tensor
record returns under $\pi$. These distinct record periods are computed
in \path{model/openwave_spin_frame_records}. They realize the spinorial
double-cover structure and do not prove an exchange-statistics theorem
or a fundamental nematic field.

\paragraph{Particle--antiparticle conversion to radiation.}
The actual native CAR current converts a neutral particle--hole pair
into two transverse quantum modes. A positive full-Hamiltonian
probability bound, complete observations and conservation of total
energy, charge and momentum establish this finite conversion mechanism.
The source uses a declared resonance, reference quench and field cap;
measurement work is retained. The calculation in
\path{model/openwave_radiative_pair} supplies no physical positronium
branching ratio or infinite-volume annihilation rate.

\paragraph{A stable neutral candidate.}
The periodic link-penalty-two source contains an isolated stationary
neutral band. Its finite polynomial/polar preparation has success above
$0.99999999999639$ and stationarity error below $4.000002\times10^{-16}$.
This supplies a stable neutral sector with complete failure records.
Exact location conservation makes its momentum band flat, so these
properties do not identify a moving relativistic dark-matter particle
or a cosmological abundance. The preparation and obstruction are both
retained in \path{model/openwave_neutral_preparation} and
\path{model/openwave_neutral_dispersion}.

\paragraph{Electrostatic spatial dependence.}
The declared Gauss Green operator gives the Coulomb $1/r$ potential and
the associated static inverse-square force in its controlled spatial
domain. A separate autonomous eighteen-state native-current experiment
supplies actual force and displacement records. This connects a
constraint law to finite dynamics; its coupling is supplied and is not
a prediction of the observed fine-structure constant. The source and
reference domains in \path{model/openwave_pinned_coulomb} and
\path{model/openwave_autonomous_coulomb} are retained.

\paragraph{Two transverse Maxwell waves.}
A finite cochain field source has two transverse Maxwell modes, exact
finite-occupation dynamics and independently computed real-space
propagation. Discrete constraints remove longitudinal radiation in the
stated sector. This reproduces the transverse wave structure of
electromagnetism on the declared lattice. A shared physical particle,
coupling and elapsed-time calibration is required for a common
electromagnetic interpretation; see
\path{model/openwave_field_observables}.

\paragraph{A massive second-order wave law.}
For the full native kinetic/pairing source, charged frames obey
\[
 h_k^2=\left(\sum_i\frac{\sin^2(ap_i)}{a^2}+e_k^2\right)I.
\]
Complete spatial wave records and a controlled low-momentum limit give
the massive Klein--Gordon form. Every finite-derivative valley and
spectator is retained. This is a source-derived wave identity with
stated momentum/spacing error, rather than an identification of all
interacting particle masses or an elimination of lattice multiplicities.

\paragraph{Discrete atomic levels.}
The declared Gauss atom has negative standing-wave levels under
independent volume and spacing variation. In its atomic units, the
first excitation gap approaches $3/8$, and the tested second-shell
splitting decreases to $0.000410646363$. This is the hydrogen-like
orbital pattern generated by the specified kinetic and Gauss source.
Mass and coupling are supplied. Precision physical hydrogen, fine
structure and radiative shifts are outside this calculation, whose
refinement evidence is in \path{model/openwave_atomic_refinement}.

\paragraph{Charged-lepton numerical comparison.}
The separate analyzer-family comparison gives electron-anchored values
$0.00051099895069$, $0.1056594144959$ and
$1.7769849708134\,\mathrm{GeV}$. The muon residual is approximately
$+9.83$ parts per million. The phase $2/9$ has retrospective selection
ancestry, and the electron fixes the common unit. These small residuals
are useful quantitative comparisons, with those empirical dependencies
retained; they are not three independent parameter-free mass predictions.
The full physical particle dictionary and interacting mass-shell
identification require further same-source tests.
